\begin{figure*}[t]
    \centering
    \includegraphics[width=\linewidth]{figure/fig2_v1.pdf}
    \caption{The overview of consistency filtering process.
    The purpose of this process is to remove misaligned data that may introduce model collapse during training. Specifically, rendered multi-view images are input into BLIP to generate captions. Then, the captions are summarized to one caption by GPT-4. Finally, the quality of the generated data is evaluated by comparing the text used for generation with the generated captions through two matching methods: word-level matching and concept-level matching.
    % Based on the generated caption, a score is calculated based on text-based and semantics concepts, and 3D models with scores below a certain level are excluded as noise data.
    }
    \label{fig:enter-label_fig4}
    \vspace{-10pt}
\end{figure*}

\section{Proposed Method: TeGA}
In this section, we introduce TeGA, expanding language-image-3D datasets with high-quality synthetic 3D data generated from text-to-3D models. 


\noindent{\textbf{Overview of TeGA.}}
To tackle the problem of 3D data scarcity in zero-shot 3D classification, the proposed method automatically generates a synthetic dataset composing language, image, and 3D modalities using a text-to-3D model.
In detail, to compose an expansion dataset, the proposed method combines the desired text, the point cloud generated by the text-to-3D model, and the images rendered from the generated point cloud.
The proposed method can generate large amounts of data without requiring human data collection or annotation and expand real datasets.
However, the generation process may err in some cases due to failure to generate or render, meaning that the alignment between each modality may not be accurate. This error may cause the model to collapse during training. 
To address this problem, we introduce a consistency filtering strategy of TeGA and the quality of the expanded dataset is improved by filtering low-quality data.
% The proposed TeGA extends the training paradigm for language-image-3D learning by going beyond data augmentation. 
% Some of the existing CLIP-based language-image-3D learning methods, which achieved SoTA at zero-shot 3D classification, have attempted to distill knowledge from CLIP with training by limited language-image-3D datasets.
% However, the lack of datasets prevents sufficient knowledge distillation from CLIP.
% % this scarcity fundamentally limits the scope of open-world 3D understanding. 
% By constructing a synthetic language-image-3D dataset from text-to-3D models, TeGA effectively compensates for the lack of datasets.
 
\noindent{\textbf{Synthetic Dataset Construction.}}
Our goal is to generate a synthetic dataset $D_{t} = \{(x_{i}^{I}, x0_{i}^{T \prime}, x_{i}^{P \prime})\}_{i=1}^{n_{o}} $ consisting of text $x_{i}^{T}$, images $x_{i}^{I \prime}$, and point clouds $x_{i}^{P \prime}$ using the text $x_{i}^{T}$ from the real dataset $D_{s} = \{(x_{i}^{I}, x_{i}^{T}, x_{i}^{P})\}_{i=1}^{n_{s}}$ that also comprises text $x_{i}^{T}$, images $x_{i}^{I}$, and point clouds $x_{i}^{P}$.
We use Point-E~\cite{nichol2022point} as the text-to-3D model of TeGA because Point-E is suitable for generating large amounts of data due to its ability to rapidly generate point clouds.
Also, we utilized only ShapeNet as the real dataset before applying TeGA instead of several datasets because our goal is not to outperform state-of-the-art methods but to serve as a training dataset expansion.
ShapeNet is a dataset containing 52,470 point cloud samples across 55 classes.
% Here, we describe the procedure for augmenting the real dataset with a synthetic dataset using TeGA, which employs Point-E as the text-to-3D model.

% In the TeGA data generation process, a synthetic 3D dataset is constructed efficiently using Point-E. ShapeNet serves as our baseline for generation cost and structure, guiding the synthetic 3D data generated through Point-E by providing ShapeNet and its category names as text prompts. The resulting synthetic dataset, termed Point-E ShapeNet, mirrors ShapeNet's structure with 52,470 samples across 55 categories. 
% Constructing Point-E ShapeNet requires approximately 55 hours using 5 NVIDIA TITAN RTX GPUs, demonstrating TeGA's scalability for practical synthetic dataset creation. \tbd{ToDo: please write the projection of synthetic images. (Torimi)}

We first generate a point cloud by inputting the category names $x_{i}^{T}$ of the real dataset $D_s$ as text prompts into Point-E.
The generation process is represented as follows:
\begin{equation}
    x_{i}^{P \prime} = G_\theta \left( x_{i}^{T}, \omega \right)
\end{equation}
% The generated point cloud is normalized, and a mesh is generated. 
We have the point clouds and text, so we render images from the point clouds. However, without mesh information, the rendered images may appear to lack accurate shape details. To address this, we perform meshing using the Ball Pivoting Algorithm~\cite{bernardini1999ball}.
While various meshing techniques exist, Point-E's point clouds are not designed for meshing, and, in our experience, using state-of-the-art meshing methods often leads to the collapse of the rendered 3D data. Therefore, we adopted the older meshing technique, the Ball Pivoting Algorithm.
Then, images $x_{i}^{I \prime}$ are generated by rendering the 3D data from 20 viewpoints.
The viewpoints are determined by rotating it clockwise in 18-degree increments, with focal lengths dynamically determined by the Open3D library.
The processes of meshing and rendering are represented as follows:
\begin{equation}
    x_{i}^{I \prime}= \mathcal{R}(\mathcal{M}(x_{i}^{P \prime}))
\end{equation}
where $\mathcal{M}(\cdot)$ denotes meshing process and $\mathcal{R}(\cdot)$ denotes rendering process.
% Then, we render the 3D model from 20 viewpoints by rotating it clockwise in 18-degree increments with focal lengths dynamically determined by the Open3D library.
% Using the generated images and text prompts, we perform the consistency filtering explained in the next section to evaluate whether the generated data matches the class labels of the real dataset.
The sets of inputted text $x_{i}^{T}$, rendered images $x_{i}^{I \prime}$, and generated point clouds $x_{i}^{P \prime}$ that pass the consistency filtering are adopted as the synthetic dataset $D_t$.

Figure~\ref{fig:enter-label_fig2} shows examples of synthetic and real 3D data generated by Point-E. The real 3D data is a sample from Objaverse-LVIS. Figure~\ref{fig:enter-label_fig2} shows that the synthetic 3D data generated by Point-E outputs a 3D shape that matches the input text. On the other hand, there are cases where even the smallest details of the 3D shape are not generated accurately. For example, the synthetic 3D data generated by Point-E has difficulty in generating even the spokes of a bicycle.
% At the rendering phase, 
% To create a dataset usable for pretraining a Language-Image-3D model, the input text for a text-to-3D model and the output point clouds are processed to align the three modalities: Language, Image, and 3D. The Language-Image-3D model learns complementarily using alignment across these modalities. Consequently, a strong alignment is required for datasets used with this model. However, point cloud data generated by text-to-3D models present the following challenges:
% (i) point cloud data must be rendered into images to generate the Image modality
% (ii) the output point clouds often do not align with the intent of the input text.

% The issue of rendering dependency can be resolved by constructing meshes. Specifically, the generated point cloud data is normalized, and meshes are created using the Ball Pivoting Algorithm~\cite{bernardini1999ball}, which generates a mesh from the point cloud information. The resulting mesh is then rendered to produce images.
% The issue of misalignment between the input text and the output point cloud is addressed using a process called Consistency Filtering. This method evaluates the alignment between the text and point cloud and filters the data accordingly. The details of Consistency Filtering are explained in the next section.

% \begin{figure*}[t]
%     \centering
%     \includegraphics[width=\linewidth]{figure/proposed_model.pdf}
%     \caption{Caption}
%     \label{fig:enter-label}
% \end{figure*}


% \begin{figure}[tb]
%     \centering
%     \includegraphics[width=\linewidth]{figure/fig3_v3_otsuka.png}
%     \caption{
%     \kt{
%     % Comparison between Real Data and Synthetic Data. The first row represents Real Data, while the second row shows Synthetic Data generated by Point-E. Comparing the red boxes between the Real Data and Synthetic Data for each class reveals that the Real Data captures finer details. On the other hand, the Synthetic Data still allows for intuitive recognition of objects such as bicycles, chairs, and beds by human observers.
%     Comparison between Real Data and Synthetic Data. The upper shows real data, while the lower shows synthetic data. Comparing the two, both capture the general appearance, but real data represents details more accurately.
%     }}
%     \label{fig:enter-label}
%     \vspace{-10pt}
% \end{figure}

\noindent\textbf{Consistency Filtering.}
Alignment across modalities is crucial for language-image-3D pretraining. 
In our generation process, alignment issues may occur at two stages: generating point clouds from text using Point-E and rendering images from point clouds. Failures in point cloud generation often result from Point-E's misinterpretation of the text. Also, rendering failures occur when the point cloud is unsuitable for meshing, such as when parts of the point cloud are incomplete. 
Training a model with such misaligned data could lead to a breakdown of modality alignment within the model. To address this, we apply consistency filtering to each generated data sample.

% Our proposed TeGA provides consistency filtering to ensure consistency between synthetic 3D data generated from text-to-3D models and text prompts. 
% The text prompts entered into the text-to-3D models in this paper are the category names of the real dataset. Therefore, if the input prompts and the generated 3D models do not match, it is possible that these are synthetic 3D models that were not included in the categories entered in the text prompts. Synthetic 3D models that do not match the text prompts may act as noise data in the training process under the data generation conditions of this paper. To address this issue, we built a filtering strategy into TeGA that excludes synthetic 3D models that are not semantically consistent with the prompts.

If alignment is maintained between the input text $x_{i}^{T}$ and the final output image $x_{i}^{I \prime}$ in the data generation process, it can be inferred that the intermediate output, the point cloud $x_{i}^{P \prime}$, is also aligned. This ensures that the alignment of the entire generated dataset is verified. Therefore, we evaluate the alignment between the text and images. Based on existing evaluation methods for 3D data~\cite{he2023t}, our filtering compares the integrated text generated from captions of multi-view images with the input text using GPT-4 and algorithm-based approaches. The process is shown in Figure~\ref{fig:enter-label_fig4}.

% Our filtering process uses two different metrics as shown in Figure~\ref{fig:enter-label_fig4}: a text-based consistency metric and a semantic alignment metric. 
For each synthetic data, we first select two images showing the front and back of the 3D data among the generated images during the dataset creation process. 
Next, we input the two images into BLIP~\cite{li2023blip} to generate each caption.
% and then we use GPT-4 to merge the two captions into a single common description.
To simplify, the generated captions for all viewpoints are combined into one unified text  $y_i$ using GPT-4.
The combined caption is then compared with the text used for generation and evaluates whether the synthetic data is aligned or not. The evaluation is conducted using two metrics: a text-based consistency metric and a semantic alignment metric. 
For the text-based consistency metric, we then compute a match score between this shared caption and the original text query. 
Specifically, if the generated caption $y_i$ contains the text $x_{i}^{T}$, a score of 5 is assigned; if it does not, a score of 1 is assigned.
\begin{equation}
\label{eq:text-based-consistency}
\text{s}_{\text{text}}(y_i, x_{i}^{T})
=
\begin{cases}
5, & \text{if } x_{i}^{T} \subseteq y_i, \\
1, & \text{otherwise}.
\end{cases}
\end{equation}
For the semantic alignment metric, inspired by the alignment evaluation in T3Bench~\cite{he2023t}, we use GPT-4 to evaluate the semantic alignment of the generated caption relative to the original prompt. 
Specifically, the common captions and text prompts generated from the two images are used to generate a five-point semantic-based similarity score using GPT-4. 
\begin{equation}
\label{eq:semantic-alignment}
\text{s}_{\text{sem}}\bigl(y_i, x_{i}^{T}\bigr)
=
\mathrm{GPT4}\bigl(y_i, x_{i}^{T}\bigr)
\;\in\;\{1, 2, 3, 4, 5\}.
\end{equation}
For more information about the GPT-4 prompts, see Appendix~\ref{app:details_tega}.

The two calculated scores $\text{s}_{\text{text}}, \text{s}_{\text{sem}}$ are summed to obtain the final score $s$. 
\begin{equation}
\label{eq:semantic-alignment}
\text{s}\bigl(y_i, x_{i}^{T}\bigr)
=
\text{s}_{\text{text}}\bigl(y_i, x_{i}^{T}\bigr) + \text{s}_{\text{sem}}\bigl(y_i, x_{i}^{T}\bigr)
\end{equation}
If this score $\text{s}\bigl(y_i, x_{i}^{T}\bigr)$ exceeds the threshold $\delta$, the data is considered aligned and included in the dataset; otherwise, it is discarded. 
We experimentally validated this threshold $\delta$ during data generation with ShapeNet and finally set it to 3.5 for use in the dataset. Filtering with this threshold in ShapeNet resulted in approximately half of the data being filtered out.
Figure~\ref{fig:filtering_result} shows samples of our consistency filtering.
The filtered data has well-formed meshes, making it visually easy to identify the objects. In contrast, the data that failed the filtering process often has collapsed meshes or represents unrecognizable objects.

% Our filtering process uses two different metrics: a text-based consistency metric and a semantic alignment metric. For each 3D model, we generate two images by projecting the 3D model from two different points of view with focal lengths that are dynamically determined by the Open3D library.
% For the text-based consistency metric, we input the two images into BLIP to generate captions, then use GPT-4 to merge these captions into a single unified description. We then compute a matching score between this unified caption and the original text prompt, evaluating consistency on a 5-point scale. The scoring method is defined as follows. For the semantic alignment metric, inspired by the alignment assessment in T3Bench, we use GPT-4 to assess the semantic coherence of the generated caption relative to the original prompt. Each of the two images is captioned via BLIP, and GPT-4 is prompted to score their alignment with the prompt on a 5-point scale. Both metrics are scored on a 5-point scale, and for this study, we set an empirical threshold of 3.5. Any 3D model scoring below this threshold on either metric is excluded as noise.